% !TeX root = ../main.tex

\chapter{Dirichlet Convolution}

\section{Composition Acting on the Coefficients}
% =================================================

Now consider

\[
F(x)
=
\sum_{n\geq1}c_n\phi_n(x).
\]

Applying \(T_m\), we obtain

\[
T_mF(x)
=
F(x^m).
\]

Therefore

\[
F(x^m)
=
\sum_{n\geq1}c_n\phi_n(x^m).
\]

Using

\[
\phi_n(x^m)=\phi_{mn}(x),
\]

we obtain

\[
\boxed{
T_mF
=
\sum_{n\geq1}c_n\phi_{mn}.
}
\]

For \(m=2\),

\[
T_2F
=
c_1\phi_2
+
c_2\phi_4
+
c_3\phi_6
+
c_4\phi_8
+\cdots.
\]

Thus the coefficient positions

\[
1,2,3,4,\ldots
\]

are sent to

\[
2,4,6,8,\ldots.
\]

For \(m=3\), they are sent to

\[
3,6,9,12,\ldots.
\]

More generally,

\[
\boxed{
n\longmapsto mn.
}
\]

% =================================================
\section{Dirichlet Convolution Appears}
% =================================================

A stronger connection appears when we combine several composition
operators.

Let

\[
A
=
\sum_{m\geq1}d_mT_m,
\]

where \(d_m\) is another sequence.

Apply \(A\) to

\[
F(x)
=
\sum_{n\geq1}c_n\phi_n(x).
\]

Then

\[
AF
=
\sum_{m\geq1}
d_mT_m
\left(
\sum_{n\geq1}c_n\phi_n
\right).
\]

Formally,

\[
AF
=
\sum_{m,n\geq1}
d_mc_n\phi_{mn}.
\]

Now collect all terms whose index is equal to \(k\).

The equation

\[
mn=k
\]

is equivalent to

\[
m\mid k,
\qquad
n=\frac{k}{m}.
\]

Therefore the coefficient of \(\phi_k\) becomes

\[
\sum_{m\mid k}
d_m c_{k/m}.
\]

Thus

\[
\boxed{
AF
=
\sum_{k\geq1}
\left(
\sum_{m\mid k}
d_m c_{k/m}
\right)
\phi_k.
}
\]

But the expression

\[
(d*c)(k)
=
\sum_{m\mid k}
d_m c_{k/m}
\]

is precisely the \emph{Dirichlet convolution} of the sequences
\(d\) and \(c\).

Therefore

\[
\boxed{
\left(
\sum_{m\geq1}d_mT_m
\right)
F
\quad\longleftrightarrow\quad
d*c.
}
\]

This gives a direct connection between composition of functions and
one of the central operations of multiplicative number theory.

\subsection*{Relation to Dirichlet convolution}

Two distinct structures coexist in the monomial‑exponential framework.

\paragraph{1. Functional layer.}
The substitution operators


\[
T_m:F(x)\mapsto F(x^m)
\]


act on functions. When these operators act on the basic blocks
\(\phi_n(x)=e^{\lambda x^n}\) we have


\[
T_m\phi_n=\phi_{mn},
\]


so composition of power maps corresponds to multiplication of indices:


\[
T_mT_n=T_{mn}.
\]



\paragraph{2. Arithmetic layer.}
If


\[
F(x)=\sum_{n\ge1} c_n\phi_n(x),
\]


and one forms a linear combination of composition operators


\[
A=\sum_{m\ge1} d_m T_m,
\]


then the action of \(A\) on \(F\) rearranges coefficients by multiplicative index relations:


\[
AF=\sum_{m,n\ge1} d_m c_n \phi_{mn}
=\sum_{k\ge1}\Big(\sum_{m\mid k} d_m c_{k/m}\Big)\phi_k.
\]


The inner sum \(\sum_{m\mid k} d_m c_{k/m}\) is precisely the Dirichlet convolution
\((d*c)(k)\). Thus the operator \(A\) corresponds to the arithmetic operation \(d*c\)
on the coefficient sequences.

\paragraph{Conclusion.}
The analytic composition operators encode multiplicative structure of indices,
while Dirichlet convolution governs how coefficient sequences combine under those operators.
Consequently, many algebraic and inversion phenomena (e.g. divisor‑sum operators and Möbius inversion)
appear naturally in this setting.


% =================================================
\section{The Divisor-Sum Operator}
% =================================================

Consider the special case

\[
d_m=1
\]

for every \(m\).

Define formally

\[
D
=
\sum_{m\geq1}T_m.
\]

Then the coefficient of \(\phi_k\) in \(DF\) is

\[
\sum_{m\mid k}c_{k/m}.
\]

Renaming the divisor,

\[
\boxed{
\sum_{d\mid k}c_d.
}
\]

Thus the operator \(D\) transforms the coefficient sequence

\[
c_k
\]

into its divisor-sum sequence

\[
\boxed{
g_k=\sum_{d\mid k}c_d.
}
\]

This is exactly the type of transformation for which Möbius inversion
is designed.

\subsection*{The Divisor‑Sum Operator: practical examples}

Define the divisor‑sum operator


\[
D=\sum_{m\ge1}T_m,\qquad T_mF(x)=F(x^m).
\]


If


\[
F(x)=\sum_{n\ge1}c_n\phi_n(x),
\]


then


\[
DF(x)=\sum_{k\ge1}\Big(\sum_{d\mid k}c_d\Big)\phi_k(x).
\]


Thus the coefficient transform is


\[
g_k=\sum_{d\mid k}c_d,
\]


i.e. \(g=\mathbf{1}*c\) (Dirichlet convolution with the constant function \(1\)).

\paragraph{Numerical example.}
Let \(c_1=1,\ c_2=2,\ c_3=0,\ c_4=1\). Then


\[
\begin{aligned}
g_1&=c_1=1,\\
g_2&=c_1+c_2=3,\\
g_3&=c_1+c_3=1,\\
g_4&=c_1+c_2+c_4=4.
\end{aligned}
\]


Hence


\[
DF(x)=\phi_1+3\phi_2+\phi_3+4\phi_4+\cdots.
\]


Recover \(c\) from \(g\) by Möbius inversion:


\[
c_n=\sum_{d\mid n}\mu(d)\,g_{n/d},
\]


which reproduces \(c_1=1,\ c_2=2,\ c_3=0,\ c_4=1\).

\paragraph{Matrix form (truncated).}
For indices \(1\le n,k\le4\) the action of \(D\) is \(g=Mc\) with


\[
M_{k,n}=\begin{cases}1,& n\mid k,\\0,&\text{otherwise,}\end{cases}
\qquad
M=\begin{pmatrix}
1 & 0 & 0 & 0\\
1 & 1 & 0 & 0\\
1 & 0 & 1 & 0\\
1 & 1 & 0 & 1
\end{pmatrix}.
\]



\paragraph{Remarks.}
The divisor‑sum operator is a concrete bridge between the operator
language \(T_m\) and multiplicative number theory: linear combinations
of the \(T_m\) correspond to Dirichlet convolutions on coefficient
sequences, and Möbius inversion provides the natural inverse procedure.


% =================================================
\section{Möbius Inversion}
% =================================================

Suppose

\[
g_n
=
\sum_{d\mid n}c_d.
\]

Then Möbius inversion gives

\[
\boxed{
c_n
=
\sum_{d\mid n}
\mu(d)
g_{n/d},
}
\]

where \(\mu\) is the Möbius function.

Therefore a classical number-theoretic inversion formula appears
naturally from the composition operators of our function family.

Symbolically,

\[
\mathbf{1}*c=g,
\]

and

\[
\mu*g=c,
\]

because

\[
\mu*\mathbf{1}=\varepsilon.
\]

Here \(\varepsilon\) is the identity for Dirichlet convolution.

This suggests a possible bridge between monomial exponential series
and arithmetic functions.

% =================================================
